% Generated draft: requires mathematical and source review.
\documentclass[11pt]{article}
\usepackage[margin=25mm]{geometry}
\usepackage{fontspec}
\setmainfont{Latin Modern Roman}
\usepackage{amsmath,amssymb,mathtools,unicode-math}
\setmathfont{Latin Modern Math}
\usepackage{enumitem,needspace,xurl,hyperref}
\setlength{\parindent}{0pt}
\setlength{\parskip}{5pt}
\setlength{\emergencystretch}{3em}
\newcommand{\N}{\mathbb{N}}
\newcommand{\Z}{\mathbb{Z}}
\newcommand{\Q}{\mathbb{Q}}
\newcommand{\R}{\mathbb{R}}
\newcommand{\C}{\mathbb{C}}
\newcommand{\F}{\mathbb{F}}
\newcommand{\A}{\mathbb{A}}
\newcommand{\PP}{\mathbb{P}}
\newcommand{\E}{\mathbb{E}}
\newcommand{\eps}{\varepsilon}
\newcommand{\dd}{\,\mathrm{d}}
\begin{document}
\section*{UnsolvedMath ID: 701534}
\textbf{Source problem code:} catalogue entry
\subsection{Definitions and mathematical statement}
Let \(G\) be a finite group and let \(\rho\colon R\to\operatorname{Aut}(G)\) be an action of a group \(R\). Define
\[
C_G(R)=\{g\in G:\rho(r)(g)=g\text{ for every }r\in R\}.
\]
For a subgroup \(H\leq G\), write
\[
C_G(H)=\{g\in G:gh=hg\text{ for every }h\in H\}.
\]
For a prime \(p\), a \(p\)-subgroup is a subgroup of order \(p^a\) for an integer \(a\geq0\), and \(O_p(G)\) denotes the largest normal \(p\)-subgroup of \(G\).

A subgroup \(Q\leq G\) is \(RC_G(R)\)-invariant if
\[
\rho(r)(Q)=Q\quad(r\in R),\qquad cQc^{-1}=Q\quad(c\in C_G(R)).
\]
Let \(P\) be the unique maximal member, with respect to inclusion, of the set of \(RC_G(R)\)-invariant \(p\)-subgroups of \(G\).

The source assumes that \(R\) acts coprimely on \(G\). Under the usual finite-group convention, this means that \(R\) is finite and \(\gcd(|R|,|G|)=1\). The supplied statement does not expressly require \(R\) to be finite or specify a convention for infinite acting groups. The formulation below retains the source's coprime-action hypothesis, with this domain qualification unresolved.

For each \(p\in\{2,3\}\), is it true that every such coprime action satisfying
\[
C_G(O_p(G))\leq O_p(G)
\]
has a subgroup \(K\) such that
\[
1\neq K\leq P,\qquad
\alpha(K)=K\ \text{for every }\alpha\in\operatorname{Aut}(P),\qquad
 gKg^{-1}=K\ \text{for every }g\in G?
\]
Thus \(K\) must be characteristic in \(P\) and normal in \(G\).

The source attributes the corresponding assertion for \(p>3\) to P. Flavell, \emph{J. Algebra} \textbf{257} (2002), 249--264. This is supplied background, not an independently verified theorem statement. The source does not exclude \(G=1\); that case satisfies the displayed centralizer hypothesis but cannot have the requested nontrivial subgroup. Consequently, the intended nontriviality restriction requires clarification for both the question and its quoted background assertion. No additional restriction has been silently imposed.
\subsection{Short English statement}
For each prime \(p=2\) and \(p=3\), ask whether the unique maximal \(p\)-subgroup invariant under a coprime action and conjugation by its fixed-point subgroup contains a nontrivial characteristic subgroup normal in \(G\), assuming \(C_G(O_p(G))\leq O_p(G)\). The acting-group convention and the omitted nontriviality restriction require clarification.
\subsection{Sources}
\begin{itemize}
\item Local UnsolvedMath registry, permanent ID \(701534\), exact source code \texttt{catalogue entry}, title \emph{Kourovka Notebook Problem 18.106}. The supplied record labels the problem open; this label has not been independently verified.
\item \emph{The Kourovka Notebook (Unsolved Problems in Group Theory)}, Problem 18.106, proposed by P. Flavell in 2014, 18th issue (2014). The record also identifies the 21st issue edition (2026), \texttt{arXiv:1401.0300v48}. The supplied collection description names E. I. Khukhro and V. D. Mazurov as editors.
\item P. Flavell, \emph{J. Algebra} \textbf{257} (2002), 249--264, cited in the supplied statement for the result when \(p>3\). No article title or direct link is supplied, and the article has not been independently consulted.
\end{itemize}
\end{document}
